% ATTEMPT_STATUS: unresolved
% ATTEMPT_ORIGIN: new research, 2026-10-01
% TOP_PROBLEM: {"id": 52, "title": "Smale's 9th Problem: Linear Programming in Polynomial Time", "category_id": 15, "category": {"id": 15, "name": "computer_science", "display_name": "Computer Science", "description": "Computational complexity, algorithms, and theoretical CS.", "slug": "computer-science", "order_index": 15, "created_at": "2026-07-31T15:26:25.671Z"}, "status": "open"}
% FIRST_POSTED: 2026-10-01
% Imported from: unsolvedmath_additional_500.tex; UnsolvedMath ID 52
% !TeX program = lualatex
% Compile twice: lualatex 52.tex
% Self-contained: no external bibliography, images, or \input files.
% Numeric section labels and visible numbers are original UnsolvedMath IDs.
\documentclass[11pt,a4paper]{article}
\usepackage[margin=24mm,headheight=14pt]{geometry}
\usepackage{fontspec}
\setmainfont{Latin Modern Roman}
\setsansfont{Latin Modern Sans}
\setmonofont{Latin Modern Mono}
\usepackage{amsmath,amssymb,mathtools}
\usepackage{unicode-math}
\setmathfont{Latin Modern Math}
\usepackage{microtype}
\usepackage{enumitem,xurl,needspace}
\usepackage[dvipsnames]{xcolor}
\usepackage{hyperref}
\hypersetup{unicode=true,colorlinks=true,linkcolor=MidnightBlue,urlcolor=MidnightBlue,
 pdftitle={UnsolvedMath Problem 52},pdfauthor={Source-annotated compilation},bookmarksnumbered=true}
\usepackage{bookmark,tocloft,titlesec,fancyhdr}
\setlength{\cftsecnumwidth}{4.2em}
\cftsetpnumwidth{2.5em}
\cftsetrmarg{4em}
\titleformat{\section}{\Large\bfseries\raggedright}{\thesection}{0.7em}{}
\pagestyle{fancy}\fancyhf{}
\fancyhead[L]{\small\sffamily UnsolvedMath research attempt}
\fancyhead[R]{\small Original catalogue IDs}
\fancyfoot[C]{\thepage}
\setlength{\parindent}{0pt}
\setlength{\parskip}{5pt plus 1pt}
\setlength{\emergencystretch}{3em}
\setcounter{tocdepth}{1}\setcounter{secnumdepth}{1}
\setlist[itemize]{leftmargin=3em,itemsep=4pt,topsep=4pt,parsep=0pt}
\allowdisplaybreaks
\newcommand{\N}{\mathbb{N}}
\newcommand{\Z}{\mathbb{Z}}
\newcommand{\Q}{\mathbb{Q}}
\newcommand{\R}{\mathbb{R}}
\newcommand{\C}{\mathbb{C}}
\newcommand{\F}{\mathbb{F}}
\newcommand{\A}{\mathbb{A}}
\newcommand{\PP}{\mathbb{P}}
\newcommand{\E}{\mathbb{E}}
\newcommand{\eps}{\varepsilon}
\newcommand{\dd}{\,\mathrm d}
\newcommand{\sref}[2]{\hyperlink{source:#1:#2}{[S#2]}}
\newif\ifshowcatalogue
\showcataloguetrue
\newcommand{\cataloguescope}[1]{\ifshowcatalogue
 \paragraph{Statement recorded in the supplied source.}
 #1\fi}
\begin{document}
% UnsolvedMath ID 52; source catalogue code SMA-009

\setcounter{section}{51}

\section{Smale's Strongly Polynomial Linear Programming Problem}\label{52}

{\small\sffamily UnsolvedMath ID 52 · Computer Science}

\subsection{Definitions and mathematical statement}

\paragraph{Shared notation.}Unless a section says otherwise, $\N=\{1,2,\ldots\}$,
$\Z$ is the ring of integers, $\Q,\R,\C$ are the rational, real, and complex
fields, $[N]=\{1,\ldots,\lfloor N\rfloor\}$, and $\log$ is the natural
logarithm. For real functions of a parameter tending to infinity,
$f\ll g$ and $f=O(g)$ mean $|f|\leq Cg$ eventually for some fixed $C$;
$f\gg g$ reverses this comparison; $f=o(g)$ means $f/g\to0$;
$f\sim g$ means $f/g\to1$; and $f\asymp g$ means both order bounds.
A subscript on $O$ or $\ll$ specifies allowed constant dependence.
The expression $n^{o(1)}$ in an upper bound means growth smaller than
$n^\epsilon$ for each fixed $\epsilon>0$. Local definitions govern any
exception, finite-field convention, or use of the same symbol for another
object. An asymptotic claim is not automatically an effective algorithm.



Consider rational linear programmes $\min\{c^Tx:Ax\leq b\}$ with $n$ variables and $m$ inequalities. Strong polynomiality means that the number of arithmetic operations and comparisons is bounded by a polynomial in $m+n$, independently of the coefficient bit lengths, while the intermediate rational numbers have encoding lengths polynomial in the input encoding length. The algorithm must also correctly detect infeasibility or an unbounded objective. \sref{52}{1} \sref{52}{2}

\paragraph{Statement recorded in the supplied source.}
Find a strongly polynomial algorithm for linear programming. \sref{52}{1}

\subsection{Short English statement}

Can every rational linear programme be solved with an arithmetic-operation count depending polynomially only on its numbers of variables and constraints?

\subsection{Research attempt}
\textbf{Status: UNRESOLVED.} We give a strongly polynomial algorithm for the one-variable case, extend it to systems with independent coordinates, and identify the coefficient-independent counting obstruction in vertex enumeration and elimination. The argument does not solve general linear programming.

\paragraph{Formulation and source check.}
Smale's original problem list [E1], consulted 1 October 2026, asks for a polynomial algorithm over the reals deciding linear-inequality feasibility. The canonical definition here additionally specifies rational inputs, intermediate bit lengths, optimization, and detection of unboundedness. All claims below use that explicit definition. A running time polynomial in the total input bit length alone does not answer it. Conversely, arithmetic-operation counts alone are insufficient if the rational numbers created have uncontrolled encodings.

\paragraph{A complete algorithm in one variable.}
Write the constraints as $a_ix\le b_i$, with rational $a_i,b_i$. If $a_i=0$, reject when $b_i<0$ and otherwise discard the row. If $a_i>0$, it gives an upper bound $x\le b_i/a_i$; if $a_i<0$, it gives a lower bound $x\ge b_i/a_i$. Let
\[
L=\max_{a_i<0} b_i/a_i,\qquad U=\min_{a_i>0} b_i/a_i,
\]
where a missing lower bound is $-\infty$ and a missing upper bound is $+\infty$. The feasible set is exactly $[L,U]$, interpreted with the relevant missing endpoints. It is empty precisely when $L>U$ or a discarded zero row was inconsistent.

For objective $cx$, if $c>0$ the minimum is $cL$ attained at $L$ when $L$ is finite, and is unbounded below otherwise. If $c<0$, use $U$ in the same way. If $c=0$, every feasible point is optimal: choose zero if allowed, otherwise a finite endpoint. For example, an interval with no finite lower endpoint always contains either zero or its finite upper endpoint, so this selection covers one-sided and fully unbounded feasible sets.

The algorithm uses one pass, $O(m)$ arithmetic operations and comparisons, independent of numerical magnitudes. Rational division and comparison use only products and quotients of pairs of input numerators and denominators. Each retained endpoint is an original ratio, so its bit length is at most a constant times the total input bit length. Cross multiplication for comparisons has the same bound. Thus both parts of the strong-polynomial definition are satisfied for $n=1$.

\paragraph{Independent-coordinate systems.}
Suppose each row of $A$ has at most one nonzero entry. Group the inequalities by their nonzero coordinate; zero rows are treated as above. Feasibility factors into a product of intervals. Optimize $c_jx_j$ independently on each interval. If one interval is empty, the whole system is infeasible. If a feasible interval is unbounded in an improving objective direction, the full problem is unbounded below: fix feasible values in all other coordinates and move only that coordinate. Otherwise the choices above yield an optimal vector.

This uses $O(m+n)$ arithmetic operations after identifying the nonzero entries, or $O(mn)$ when reading a dense input matrix. Both counts are polynomial in $m+n$. The bit-length argument remains coordinatewise and introduces no iterative coefficient growth. This is a complete special-class theorem, not a claim about sparse matrices in general: a row involving two coordinates already couples their feasible intervals.

\paragraph{Why vertex enumeration does not finish the problem.}
For a bounded full-dimensional feasible polytope, a linear objective has an optimal vertex. Enumerating $n$-element subsets of the $m$ constraints, solving each full-rank equality system, and testing feasibility will therefore find an optimum in this restricted case. Even granting all linear-algebra and bit-length details, the number of subsets is $\binom mn$. For $m=2n$,
\[
\binom{2n}{n}\ge\frac{2^{2n}}{2n+1},
\]
because the middle binomial coefficient is the largest among $2n+1$ coefficients summing to $2^{2n}$. The enumeration count is exponential in $n$. Its independence from coefficient size does not make it polynomial in dimension.

The restriction to bounded full-dimensional polytopes is also substantive. A nonempty affine line has no vertices, while a constant objective on it has a finite attained optimum. An algorithm that reports infeasibility merely because it found no basic feasible point would fail on such an input. Adding an arbitrary large box requires a justified size bound and cannot silently repair this issue.

\paragraph{Fourier--Motzkin elimination: a second failed route.}
To eliminate one variable, solve every row with a positive coefficient for an upper bound on that variable and every row with a negative coefficient for a lower bound. Pairing each lower bound with each upper bound enforces the necessary and sufficient compatibility inequalities in the remaining variables. With $p$ lower and $q$ upper bounds, this step can create $pq$ rows, in addition to zero-coefficient rows. Sufficiency follows because any assignment satisfying every pair has maximum lower bound at most minimum upper bound, so an eliminated variable can be selected between them.

This is an exact feasibility reduction, but it is not a polynomial-size guarantee. A roughly balanced split can square the number of rows at a step; repeated elimination yields a rapidly growing a priori bound. This observation is not a lower bound for every conceivable elimination algorithm, because redundancy removal could discard many rows. The missing theorem would have to show that enough rows can always be discarded efficiently while preserving the projected feasible set. We have not proved such a theorem.

\paragraph{Adversarial verification and remaining gap.}
Zero coefficients, inconsistent constant rows, both signs of the objective, infinite endpoints, and empty intervals were explicitly handled in the one-variable proof. A chain of constraints may have extremely small or large rational endpoints without changing its arithmetic-operation count; exact rational comparisons prevent numerical tolerance errors. These checks do not transfer separability to a coupled system.

The first unresolved step is a dimension-polynomial bound on the work needed for arbitrary coupled inequalities, while maintaining polynomial intermediate bit lengths and valid infeasibility or unboundedness conclusions. The two failed routes isolate different obstructions: exponential numbers of bases and uncontrolled numbers of projected constraints. No numerical experiment, new general complexity lower bound, or solution of Smale's problem is claimed.

\subsection{Sources}

{\small

\begin{itemize}
\item[E1] S. Smale, \emph{Mathematical Problems for the Next Century} (1998), Problem 9. Original problem-list text consulted 1 October 2026. \url{https://www.fim.uni-passau.de/fileadmin/dokumente/fakultaeten/fim/lehrstuhl/muller/SmaleProblems1998.pdf}.


\item[\hypertarget{source:52:1}{S1}] ULAM.AI, \emph{UnsolvedMath}, supplied \texttt{problems.json}, record 52 (SMA-009), statement and background. The embedded literature-review date is 2026-08-17; that is the archive’s review date, not a fresh certification. Dataset landing page: \url{https://huggingface.co/datasets/ulamai/UnsolvedMath}.

\item[\hypertarget{source:52:2}{S2}] Smale's Problems, MathWorld (). \url{https://mathworld.wolfram.com/SmalesProblems.html}. Bibliographic information imported from the supplied record.

\end{itemize}

}

\label{end:52}
\end{document}
